%% evaluation-spectrum.tex — the evaluability spectrum (Chapter 7). House conceptual style.
%% Verifiable oracle (left) contrasted with a model-as-judge critic (right), with the
%% critic/evaluator convergence noted below.
\input{figures/concept-style}
\begin{tikzpicture}
% the spectrum the two poles sit on
\draw[conceptflow, <->] (-2.1,1.2) -- (8.1,1.2);
\node[conceptlabel, anchor=south] at (-1.3,1.25) {verifiable};
\node[conceptlabel, anchor=south] at (7.3,1.25) {judgment};
% the two poles
\node[conceptbox, fill=gray!15,   minimum width=3.2cm] (ev) at (0,0) {Evaluator};
\node[conceptbox, fill=violet!15, minimum width=3.2cm] (cr) at (6,0) {Critic};
% distinctions, listed below each (outside the node)
\node[conceptlabel, text width=4.6cm, anchor=north] at (0,-0.75)
  {an oracle: formula, test, or check\\objective and automatic\\unambiguous, cheap\\sphere packing, code, MSE};
\node[conceptlabel, text width=4.6cm, anchor=north] at (6,-0.75)
  {a capable model scoring a rubric\\automatic but not objective\\contextual, biased, gameable\\proof, prose, an idea};
% the convergence, along the bottom
\node[conceptlabel, text width=10.5cm, anchor=north, text=black!55] at (3,-2.75)
  {Distinct where a formula exists (the critic only filters what the oracle scores); they merge where none does (the critic \emph{is} the evaluator).};
\end{tikzpicture}
